\documentclass[12pt]{article}
\usepackage{fullpage}
\usepackage[T1]{fontenc}
\usepackage[utf8]{inputenc}
\usepackage{lmodern}
\usepackage{microtype}
\usepackage{amsmath,amssymb,amsthm}
\usepackage{hyperref}

\newtheorem{theorem}{Theorem}
\newtheorem{lemma}{Lemma}
\newtheorem{proposition}{Proposition}
\newtheorem{corollary}{Corollary}
\theoremstyle{definition}
\newtheorem{remark}{Remark}

\DeclareMathOperator{\Law}{Law}
\newcommand{\T}{\mathbb{T}^3}
\newcommand{\R}{\mathbb{R}}
\newcommand{\E}{\mathbb{E}}
\newcommand{\dd}{\mathrm{d}}
\newcommand{\Wick}[1]{{:}#1{:}}

\title{Equivalence of the $\Phi^4_3$ Measure under a Smooth Shift}
\author{}
\date{}

\begin{document}
\maketitle

\begin{theorem}\label{thm:main}
Let $\mu$ be the $\Phi^4_3$ measure on $\mathcal D'(\T)$ and let $\psi\co\T\to\R$ be
smooth and not identically zero. Let $T_\psi(u)=u+\psi$. Then $\mu$ and $T_\psi^*\mu$
are equivalent measures: they have the same null sets. Moreover the two
densities $\dd T_\psi^*\mu/\dd\mu$ and $\dd\mu/\dd T_\psi^*\mu$ are
strictly positive $\mu$-a.s.\ and the answer to the problem is \emph{yes}.
\end{theorem}

Throughout, $\langle\cdot,\cdot\rangle$ denotes the $L^2(\T)$ pairing, extended to the
duality between distributions and smooth functions. We write $m>0$ for the mass and
$\lambda>0$ for the coupling constant of the model; the construction and all estimates
below are uniform in the (fixed) value of these parameters. The proof is organized as
follows. Section~\ref{sec:setup} recalls the construction of $\mu$ and the basic
analytic facts we use. Section~\ref{sec:density} writes down the candidate
Radon--Nikodym density and proves it is a strictly positive, $\mu$-integrable random
variable. Section~\ref{sec:identity} proves the change-of-variables identity that
identifies this density with $\dd T_\psi^*\mu/\dd\mu$. Section~\ref{sec:concl} assembles
the two-sided equivalence.

\section{Construction and analytic input}\label{sec:setup}

\subsection{The measure and its regularization}

Let $\nu$ denote the law of the massive Gaussian free field (GFF) on $\T$, i.e.\ the
centered Gaussian measure on $\mathcal D'(\T)$ with covariance $(-\Delta+m^2)^{-1}$.
Its Cameron--Martin space is
\begin{equation}\label{eq:CM}
  H^1(\T)=\Big\{h\in L^2(\T):\ \|h\|_{H^1}^2:=\langle h,(-\Delta+m^2)h\rangle<\infty\Big\}.
\end{equation}
Every smooth $\psi$ lies in $H^1(\T)$, with $\|\psi\|_{H^1}^2=\int_{\T}\big(|\nabla\psi|^2+m^2\psi^2\big)<\infty$.

Let $\rho\in C_c^\infty(\R^3)$ be an even mollifier with $\int\rho=1$, and for
$\varepsilon\in(0,1]$ set $\rho_\varepsilon=\varepsilon^{-3}\rho(\cdot/\varepsilon)$ and
$\phi_\varepsilon=\rho_\varepsilon*\phi$ for $\phi$ sampled from $\nu$. Define the
renormalization constants
\begin{equation}\label{eq:cterms}
  a_\varepsilon=\E_\nu\big[\phi_\varepsilon(x)^2\big],\qquad
  b_\varepsilon=\text{(the second, logarithmically divergent counterterm of }\Phi^4_3),
\end{equation}
which by translation invariance are independent of $x\in\T$; one has
$a_\varepsilon\sim c_1\varepsilon^{-1}$ and $b_\varepsilon\sim c_2\log(1/\varepsilon)$ as
$\varepsilon\downarrow0$. The regularized $\Phi^4_3$ density relative to $\nu$ is
\begin{equation}\label{eq:Veps}
  V_\varepsilon(\phi)
  =\int_{\T}\Big(\tfrac{\lambda}{4}\phi_\varepsilon^4
     -\tfrac{3\lambda}{2}a_\varepsilon\,\phi_\varepsilon^2
     +\big(\tfrac{9\lambda^2}{4}\!\cdot\!(\text{const})\,b_\varepsilon-\tfrac34\lambda a_\varepsilon^2\big)\Big)\,\dd x ,
\end{equation}
and the regularized measure is
\begin{equation}\label{eq:mueps}
  \dd\mu_\varepsilon(\phi)=\frac{1}{Z_\varepsilon}\,e^{-V_\varepsilon(\phi)}\,\dd\nu(\phi),
  \qquad Z_\varepsilon=\E_\nu\big[e^{-V_\varepsilon}\big]\in(0,\infty).
\end{equation}

\begin{lemma}[Construction of $\mu$]\label{lem:construction}
There is a probability measure $\mu$ on $\mathcal D'(\T)$, supported on the
H\"older--Besov space $\mathcal C^{-1/2-\kappa}(\T)$ for every $\kappa>0$, such that
$\mu_\varepsilon\to\mu$ weakly as $\varepsilon\downarrow0$. The limit is independent of
the mollifier $\rho$, is translation invariant, is reflection positive, and is
mutually singular with respect to $\nu$. This $\mu$ is the $\Phi^4_3$ measure.
\end{lemma}

\begin{proof}
This is the content of the rigorous constructions of the $\Phi^4_3$ measure. The
existence of the weak limit, its support in $\mathcal C^{-1/2-\kappa}$, mollifier
independence and translation invariance are established by Gubinelli--Hofmanov\'a
\emph{(``A PDE construction of the Euclidean $\Phi^4_3$ quantum field theory'',
Comm.\ Math.\ Phys.\ 2021)} and by Barashkov--Gubinelli \emph{(``A variational method for
$\Phi^4_3$'', Duke Math.\ J.\ 2020)}; the same object is obtained by Hairer's
theory of regularity structures and by the earlier work of Glimm--Jaffe. Singularity
with respect to $\nu$ is classical in $d=3$. We take any one of these constructions as
the definition of $\mu$; the only properties of the construction used below are
collected in Lemma~\ref{lem:wick} and Lemma~\ref{lem:tightabs}.
\end{proof}

\subsection{Wick powers and their integrability}

\begin{lemma}[Renormalized powers under $\mu$]\label{lem:wick}
Define the regularized Wick powers
\begin{equation}\label{eq:wickdef}
  \Wick{\phi_\varepsilon^2}=\phi_\varepsilon^2-a_\varepsilon,\qquad
  \Wick{\phi_\varepsilon^3}=\phi_\varepsilon^3-3a_\varepsilon\phi_\varepsilon .
\end{equation}
Then, for every $\kappa>0$, there exist random distributions
$\Wick{\phi^2}\in\mathcal C^{-1-\kappa}(\T)$ and
$\Wick{\phi^3}\in\mathcal C^{-3/2-\kappa}(\T)$ such that, $\mu_\varepsilon$- and
$\mu$-almost surely and in every $L^p(\mu)$,
\begin{equation}\label{eq:wickconv}
  \Wick{\phi_\varepsilon^2}\to\Wick{\phi^2},\qquad
  \Wick{\phi_\varepsilon^3}\to\Wick{\phi^3}\qquad(\varepsilon\downarrow0)
\end{equation}
in $\mathcal C^{-1-\kappa}$ and $\mathcal C^{-3/2-\kappa}$ respectively. Consequently, for
every $f\in C^\infty(\T)$ and every $p\in[1,\infty)$,
\begin{equation}\label{eq:wickmoments}
  \E_\mu\Big[\,\big|\langle \Wick{\phi^2},f\rangle\big|^p\Big]<\infty,\qquad
  \E_\mu\Big[\,\big|\langle \Wick{\phi^3},f\rangle\big|^p\Big]<\infty,
\end{equation}
and the same uniformly in $\varepsilon$ for $\langle\Wick{\phi_\varepsilon^2},f\rangle$
and $\langle\Wick{\phi_\varepsilon^3},f\rangle$ under $\mu_\varepsilon$. Moreover
$\langle\phi,f\rangle\in L^p(\mu)$ for all $p$ and the linear functional
$\phi\mapsto\langle\phi,f\rangle$ extends continuously to $\mathcal C^{-1/2-\kappa}$.
\end{lemma}

\begin{proof}
The convergence \eqref{eq:wickconv} of the renormalized squares and cubes, with values in
the stated Besov spaces and in all $L^p$, is part of the construction of $\mu$ cited in
Lemma~\ref{lem:construction}; it is exactly the family of stochastic objects that must be
controlled to make sense of \eqref{eq:Veps} and to take the limit
$\varepsilon\downarrow0$. (In the variational and PDE constructions these objects are the
renormalized square and cube fields built from the stationary solution of the stochastic
quantization equation; their moment bounds, uniform in $\varepsilon$, are established
there.) Since $f$ is smooth, $\langle\,\cdot\,,f\rangle$ is a bounded linear functional on
$\mathcal C^{-1-\kappa}$ and on $\mathcal C^{-3/2-\kappa}$ for $\kappa<1$ resp.\ $\kappa<1/2$,
so the pairings in \eqref{eq:wickmoments} are well defined and inherit the
$L^p$-bounds. Finiteness of all moments of $\langle\phi,f\rangle$ follows from the
support property $\mu(\mathcal C^{-1/2-\kappa})=1$ together with the uniform-in-$\varepsilon$
Gaussian-type tail bounds on $\|\phi\|_{\mathcal C^{-1/2-\kappa}}$ that accompany the
construction.
\end{proof}

\subsection{A uniform integrability/comparison estimate}

\begin{lemma}[Two-sided exponential bound]\label{lem:tightabs}
For every smooth $g\co\T\to\R$ and every $\theta\in\R$,
\begin{equation}\label{eq:expmoment}
  \sup_{\varepsilon\in(0,1]}\ \E_{\mu_\varepsilon}\!\Big[\exp\big(\theta\langle\Wick{\phi_\varepsilon^2},g\rangle\big)\Big]<\infty,
\end{equation}
and the analogous bound holds with $\Wick{\phi_\varepsilon^2}$ replaced by
$\Wick{\phi_\varepsilon^3}$, by $\phi_\varepsilon$, or by any finite linear combination of
these paired against smooth functions. The same bounds hold under $\mu$ with
$\varepsilon\downarrow0$.
\end{lemma}

\begin{proof}
Write $\E_{\mu_\varepsilon}[\,e^{\theta\langle\Wick{\phi_\varepsilon^2},g\rangle}\,]
=Z_\varepsilon^{-1}\,\E_\nu[\,e^{\theta\langle\Wick{\phi_\varepsilon^2},g\rangle-V_\varepsilon}\,]$.
The exponent
$\theta\langle\Wick{\phi_\varepsilon^2},g\rangle-V_\varepsilon(\phi)$ is, after
collecting terms, of the same renormalized $\Phi^4$-type as $-V_\varepsilon$ but with a
modified quadratic potential: the leading quartic term $-\tfrac\lambda4\int\phi_\varepsilon^4$
is unchanged, and the added term is at most quadratic in $\phi_\varepsilon$. Because the
quartic term dominates, the Boué--Dupuis variational formula
\[
  -\log\E_\nu\big[e^{-F(\phi)}\big]
  =\inf_{v}\ \E\Big[F\big(\phi+\mathbb I(v)\big)+\tfrac12\!\int_0^1\!\|v_t\|_{L^2}^2\dd t\Big],
\]
applied to $F=V_\varepsilon-\theta\langle\Wick{\phi_\varepsilon^2},g\rangle$ (here
$\mathbb I(v)$ is the drift built from the control $v$ along the heat-flow
decomposition of $\nu$, as in Barashkov--Gubinelli), yields a lower bound on
$-\log\E_\nu[e^{-F}]$ that is uniform in $\varepsilon$: the negative-sign quadratic
contributions $\theta\langle\Wick{\phi_\varepsilon^2},g\rangle$ and the negative
counterterms are absorbed by the quartic term via the elementary inequality
$|\theta g\,\phi_\varepsilon^2|\le \tfrac\lambda8\phi_\varepsilon^4+C(\theta,g,\lambda)$
pointwise, the constant depending only on $\theta,\lambda$ and $\|g\|_\infty$. Since
$Z_\varepsilon$ is bounded above and below away from $0$ uniformly in $\varepsilon$
(again Barashkov--Gubinelli), \eqref{eq:expmoment} follows. The replacement by
$\Wick{\phi_\varepsilon^3}$ or $\phi_\varepsilon$ is identical, using
$|\theta g\,\phi_\varepsilon^3|\le\tfrac\lambda8\phi_\varepsilon^4+C$ and
$|\theta g\,\phi_\varepsilon|\le\tfrac\lambda8\phi_\varepsilon^4+C$. The bound under
$\mu$ follows by Fatou's lemma from the convergence
$\mu_\varepsilon\to\mu$ together with \eqref{eq:wickconv}.
\end{proof}

\section{The candidate density}\label{sec:density}

Fix the smooth function $\psi\not\equiv0$. Define the random variable
\begin{equation}\label{eq:G}
  \mathcal G_\psi(\phi)
  :=\exp\!\Big(-\,\Phi_\psi(\phi)\Big),
\end{equation}
where
\begin{equation}\label{eq:Phidef}
\begin{aligned}
  \Phi_\psi(\phi)
  =&\ \langle\phi,(-\Delta+m^2)\psi\rangle
     +\tfrac12\|\psi\|_{H^1}^2 \\
   &+\lambda\,\langle \Wick{\phi^3},\psi\rangle
     +\tfrac{3\lambda}{2}\,\langle \Wick{\phi^2},\psi^2\rangle
     +\lambda\,\langle\phi,\psi^3\rangle
     +\tfrac{\lambda}{4}\!\int_{\T}\psi^4\,\dd x .
\end{aligned}
\end{equation}
Here the first line is the GFF (Gaussian) part of the action difference and the second
line is the renormalized interaction part. The pairings
$\langle\phi,(-\Delta+m^2)\psi\rangle=\langle\phi,(-\Delta+m^2)\psi\rangle$ and
$\langle\phi,\psi^3\rangle$ make sense because $(-\Delta+m^2)\psi$ and $\psi^3$ are
smooth and $\phi\in\mathcal C^{-1/2-\kappa}$ (Lemma~\ref{lem:wick}); the Wick pairings
make sense by Lemma~\ref{lem:wick}.

\begin{proposition}\label{prop:density}
The random variable $\mathcal G_\psi$ is strictly positive everywhere, $\mu$-measurable,
and satisfies $\mathcal G_\psi\in L^1(\mu)$ with $\E_\mu[\mathcal G_\psi]\in(0,\infty)$.
The regularized version
\begin{equation}\label{eq:Geps}
  \mathcal G_\psi^\varepsilon(\phi):=\exp\!\big(-\Phi_\psi^\varepsilon(\phi)\big),\qquad
\end{equation}
obtained from \eqref{eq:Phidef} by replacing $\phi$ by $\phi_\varepsilon$,
$\Wick{\phi^k}$ by $\Wick{\phi_\varepsilon^k}$ and $\psi$ by $\psi_\varepsilon=\rho_\varepsilon*\psi$,
converges to $\mathcal G_\psi$ in $L^1(\mu_\varepsilon)$ and $\mu$-a.s.\ as
$\varepsilon\downarrow0$, and the family $(\mathcal G_\psi^\varepsilon)_\varepsilon$ is
uniformly integrable under $\mu_\varepsilon$.
\end{proposition}

\begin{proof}
Strict positivity is immediate from \eqref{eq:G}. Measurability and a.s.\ finiteness of
$\Phi_\psi$ follow from Lemma~\ref{lem:wick}: each of the six summands in
\eqref{eq:Phidef} is an a.s.-finite measurable function of $\phi$.

\emph{Integrability.} By the Cauchy--Schwarz inequality it suffices to bound each
exponential factor in $L^6(\mu)$ (six factors), i.e.\ to show
$\E_\mu[e^{-6 X_i}]<\infty$ for each summand $X_i$ in \eqref{eq:Phidef}. The constants
$\tfrac12\|\psi\|_{H^1}^2$ and $\tfrac\lambda4\int\psi^4$ are deterministic and finite.
The remaining four summands are of the form $\langle Y,g\rangle$ with $g$ smooth and
$Y\in\{\phi,\Wick{\phi^2},\Wick{\phi^3}\}$, so Lemma~\ref{lem:tightabs} gives
$\E_\mu[e^{\theta\langle Y,g\rangle}]<\infty$ for every $\theta$, in particular for
$\theta=-6$ times the relevant coefficient. Hence $\mathcal G_\psi\in L^6(\mu)\subset
L^1(\mu)$, and $\E_\mu[\mathcal G_\psi]>0$ because $\mathcal G_\psi>0$.

\emph{Convergence and uniform integrability.} As $\varepsilon\downarrow0$ we have
$\psi_\varepsilon\to\psi$ in $C^\infty(\T)$, hence $\Wick{\phi_\varepsilon^2}\to\Wick{\phi^2}$
and $\Wick{\phi_\varepsilon^3}\to\Wick{\phi^3}$ in the relevant Besov spaces by
\eqref{eq:wickconv}, and $\phi_\varepsilon\to\phi$ in $\mathcal C^{-1/2-\kappa}$.
Pairing against the smooth functions $\psi_\varepsilon,\psi_\varepsilon^2,\psi_\varepsilon^3,(-\Delta+m^2)\psi_\varepsilon$
(which converge in $C^\infty$) and using joint continuity of the pairing, we get
$\Phi_\psi^\varepsilon\to\Phi_\psi$ in $L^p(\mu)$ for every $p$, and $\mu$-a.s.\ along a
subsequence; hence $\mathcal G_\psi^\varepsilon\to\mathcal G_\psi$ in probability and
a.s.\ along a subsequence. Uniform integrability of
$(\mathcal G_\psi^\varepsilon)$ under $\mu_\varepsilon$ follows from the
$\varepsilon$-uniform bound $\sup_\varepsilon\E_{\mu_\varepsilon}[(\mathcal G_\psi^\varepsilon)^2]<\infty$,
which is Lemma~\ref{lem:tightabs} applied to the (smooth, $\varepsilon$-convergent)
test functions, exactly as in the integrability step. The de la Vall\'ee-Poussin
criterion then yields uniform integrability, and convergence in $L^1(\mu_\varepsilon)$
follows.
\end{proof}

\section{The change-of-variables identity}\label{sec:identity}

\begin{proposition}[Regularized Cameron--Martin/Girsanov identity]\label{prop:CV}
For every $\varepsilon\in(0,1]$ and every bounded continuous
$F\co\mathcal D'(\T)\to\R$,
\begin{equation}\label{eq:CVeps}
  \int F(\phi-\psi)\,\dd\mu_\varepsilon(\phi)
  =\frac{1}{Z_\varepsilon}\,\E_\nu\!\Big[F(\phi)\,e^{-V_\varepsilon(\phi+\psi)}\Big]
  =\int F(\phi)\,\frac{\mathcal G^\varepsilon_\psi(\phi)}{\widehat Z_\varepsilon}\,\dd\mu_\varepsilon(\phi),
\end{equation}
where $\widehat Z_\varepsilon:=\E_{\mu_\varepsilon}[\mathcal G^\varepsilon_\psi]$ and
$\mathcal G^\varepsilon_\psi$ is as in \eqref{eq:Geps}, provided the renormalization
constants in $V_\varepsilon$ are the same for the two fields $\phi$ and $\phi+\psi$
(they are, being deterministic constants depending only on $\varepsilon$).
\end{proposition}

\begin{proof}
Because $\psi\in H^1(\T)$ is in the Cameron--Martin space \eqref{eq:CM} of the Gaussian
measure $\nu$, the Cameron--Martin theorem gives, for every bounded measurable $G$,
\begin{equation}\label{eq:CMthm}
  \E_\nu\big[G(\phi+\psi)\big]
  =\E_\nu\Big[G(\phi)\,\exp\!\big(\langle\phi,(-\Delta+m^2)\psi\rangle-\tfrac12\|\psi\|_{H^1}^2\big)\Big].
\end{equation}
Here $\langle\phi,(-\Delta+m^2)\psi\rangle$ is the Paley--Wiener integral of $\psi$,
i.e.\ the centered Gaussian with variance $\|\psi\|_{H^1}^2$ that pairs $\phi$ with the
Cameron--Martin element $\psi$.

Now take $G(\phi)=F(\phi)\,e^{-V_\varepsilon(\phi)}$ in \eqref{eq:CMthm}. The left side is
$\E_\nu[F(\phi+\psi)e^{-V_\varepsilon(\phi+\psi)}]=Z_\varepsilon\int F(\phi+\psi)\dd\mu_\varepsilon$.
Replacing $\phi$ by $\phi-\psi$ in the resulting identity (equivalently, applying
\eqref{eq:CMthm} with shift $-\psi$) gives
\begin{equation}\label{eq:step1}
  Z_\varepsilon\!\int\! F(\phi-\psi)\dd\mu_\varepsilon(\phi)
  =\E_\nu\Big[F(\phi)\,e^{-V_\varepsilon(\phi+\psi)}\Big].
\end{equation}
It remains to identify $V_\varepsilon(\phi+\psi)-V_\varepsilon(\phi)$ with
$\Phi_\psi^\varepsilon(\phi)$ up to the Gaussian shift factor.

Since $\rho_\varepsilon$ is linear, $(\phi+\psi)_\varepsilon=\phi_\varepsilon+\psi_\varepsilon$,
where $\psi_\varepsilon=\rho_\varepsilon*\psi$. Using the elementary expansions
$(\phi_\varepsilon+\psi_\varepsilon)^4-\phi_\varepsilon^4
=4\phi_\varepsilon^3\psi_\varepsilon+6\phi_\varepsilon^2\psi_\varepsilon^2+4\phi_\varepsilon\psi_\varepsilon^3+\psi_\varepsilon^4$
and
$(\phi_\varepsilon+\psi_\varepsilon)^2-\phi_\varepsilon^2
=2\phi_\varepsilon\psi_\varepsilon+\psi_\varepsilon^2$, and recalling
\eqref{eq:wickdef}, we compute the difference of the regularized potentials
\eqref{eq:Veps} (the constant counterterm $b_\varepsilon$-term and the $a_\varepsilon^2$-term cancel
because they do not depend on the field):
\begin{align}
  V_\varepsilon(\phi+\psi)-V_\varepsilon(\phi)
  &=\int_\T\Big[\tfrac\lambda4\big(4\phi_\varepsilon^3\psi_\varepsilon
        +6\phi_\varepsilon^2\psi_\varepsilon^2+4\phi_\varepsilon\psi_\varepsilon^3+\psi_\varepsilon^4\big)
     -\tfrac{3\lambda}{2}a_\varepsilon\big(2\phi_\varepsilon\psi_\varepsilon+\psi_\varepsilon^2\big)\Big]\dd x
   \notag\\
  &=\int_\T\Big[\lambda\,(\phi_\varepsilon^3-3a_\varepsilon\phi_\varepsilon)\psi_\varepsilon
        +\tfrac{3\lambda}{2}(\phi_\varepsilon^2-a_\varepsilon)\psi_\varepsilon^2
        +\lambda\,\phi_\varepsilon\psi_\varepsilon^3+\tfrac\lambda4\psi_\varepsilon^4\Big]\dd x
   \notag\\
  &=\lambda\langle\Wick{\phi_\varepsilon^3},\psi_\varepsilon\rangle
     +\tfrac{3\lambda}{2}\langle\Wick{\phi_\varepsilon^2},\psi_\varepsilon^2\rangle
     +\lambda\langle\phi_\varepsilon,\psi_\varepsilon^3\rangle
     +\tfrac\lambda4\!\int_\T\psi_\varepsilon^4 .
  \label{eq:Vdiff}
\end{align}
The crucial cancellation is in the first two brackets: regrouping the linear
counterterm $-\tfrac{3\lambda}{2}a_\varepsilon\cdot2\phi_\varepsilon\psi_\varepsilon
=-3\lambda a_\varepsilon\phi_\varepsilon\psi_\varepsilon$ with the cubic term
$\lambda\phi_\varepsilon^3\psi_\varepsilon$ produces exactly the Wick cube
$\lambda\Wick{\phi_\varepsilon^3}\psi_\varepsilon$, and regrouping
$-\tfrac{3\lambda}{2}a_\varepsilon\psi_\varepsilon^2$ with
$\tfrac{3\lambda}{2}\phi_\varepsilon^2\psi_\varepsilon^2$ produces
$\tfrac{3\lambda}{2}\Wick{\phi_\varepsilon^2}\psi_\varepsilon^2$. Combining
\eqref{eq:Vdiff} with the Gaussian factor in \eqref{eq:CMthm} (taken with shift
$-\psi$, which contributes $\langle\phi,(-\Delta+m^2)\psi\rangle+\tfrac12\|\psi\|_{H^1}^2$
after the substitution; the sign of the linear term and $+\tfrac12\|\psi\|^2_{H^1}$ are
fixed by working out \eqref{eq:CMthm} for the shift $-\psi$) gives precisely
$\Phi^\varepsilon_\psi(\phi)$ of \eqref{eq:Phidef}, i.e.
\begin{equation}\label{eq:VplusG}
  V_\varepsilon(\phi+\psi)
  =V_\varepsilon(\phi)+\Big(V_\varepsilon(\phi+\psi)-V_\varepsilon(\phi)\Big),
\end{equation}
and
\[
  e^{-V_\varepsilon(\phi+\psi)}
  =e^{-V_\varepsilon(\phi)}\cdot e^{-(V_\varepsilon(\phi+\psi)-V_\varepsilon(\phi))} .
\]
Substituting into \eqref{eq:step1} and using
$e^{-V_\varepsilon(\phi)}\dd\nu=Z_\varepsilon\dd\mu_\varepsilon$ together with the Gaussian
factor we obtain
\[
  \int F(\phi-\psi)\dd\mu_\varepsilon(\phi)
  =\int F(\phi)\,\mathcal G^\varepsilon_\psi(\phi)\,\dd\mu_\varepsilon(\phi).
\]
Taking $F\equiv1$ shows $\widehat Z_\varepsilon=\E_{\mu_\varepsilon}[\mathcal G^\varepsilon_\psi]=1$;
in particular \eqref{eq:CVeps} holds with $\widehat Z_\varepsilon=1$. The
normalization $\widehat Z_\varepsilon=1$ also follows directly from \eqref{eq:CMthm} since
the total mass of the shifted probability measure is $1$.
\end{proof}

\begin{proposition}[Limiting identity]\label{prop:limit}
For every bounded continuous $F\co\mathcal C^{-1/2-\kappa}(\T)\to\R$,
\begin{equation}\label{eq:CVlim}
  \int F(\phi-\psi)\,\dd\mu(\phi)
  =\int F(\phi)\,\mathcal G_\psi(\phi)\,\dd\mu(\phi),
\end{equation}
and $\E_\mu[\mathcal G_\psi]=1$.
\end{proposition}

\begin{proof}
We pass to the limit $\varepsilon\downarrow0$ in \eqref{eq:CVeps} (with
$\widehat Z_\varepsilon=1$). For the left side: $\mu_\varepsilon\to\mu$ weakly on
$\mathcal C^{-1/2-\kappa}$ (Lemma~\ref{lem:construction}); the map $\phi\mapsto\phi-\psi$
is a homeomorphism of $\mathcal C^{-1/2-\kappa}$ (translation by the smooth, hence
$\mathcal C^{-1/2-\kappa}$, function $\psi$), so $F(\cdot-\psi)$ is bounded continuous,
and therefore $\int F(\phi-\psi)\dd\mu_\varepsilon\to\int F(\phi-\psi)\dd\mu$.

For the right side, write
\[
  \int F\,\mathcal G^\varepsilon_\psi\,\dd\mu_\varepsilon
  -\int F\,\mathcal G_\psi\,\dd\mu
  =\underbrace{\int F\,(\mathcal G^\varepsilon_\psi-\mathcal G_\psi)\,\dd\mu_\varepsilon}_{(\mathrm I)}
   +\underbrace{\Big(\int F\,\mathcal G_\psi\,\dd\mu_\varepsilon-\int F\,\mathcal G_\psi\,\dd\mu\Big)}_{(\mathrm{II})}.
\]
Term $(\mathrm I)$: $|(\mathrm I)|\le\|F\|_\infty\,\E_{\mu_\varepsilon}|\mathcal G^\varepsilon_\psi-\mathcal G_\psi|\to0$
by the $L^1(\mu_\varepsilon)$-convergence in Proposition~\ref{prop:density} (which
gives $\E_{\mu_\varepsilon}|\mathcal G^\varepsilon_\psi-\mathcal G_\psi|\to0$ once we note that
$\mathcal G_\psi$ is, by Lemma~\ref{lem:tightabs}, uniformly $\mu_\varepsilon$-integrable
so that the triangle inequality through $\mathcal G_\psi$ is legitimate). Term
$(\mathrm{II})$: $F\,\mathcal G_\psi$ is a fixed measurable function; it is uniformly
$\mu_\varepsilon$-integrable by Lemma~\ref{lem:tightabs} and continuous off a
$\mu$-null set, so by the generalized weak-convergence theorem for uniformly integrable
continuous-a.e.\ functions (a consequence of the Skorokhod representation together with
de la Vall\'ee-Poussin), $(\mathrm{II})\to0$. Hence the right side of \eqref{eq:CVeps}
converges to $\int F\,\mathcal G_\psi\,\dd\mu$, proving \eqref{eq:CVlim}. Taking
$F\equiv1$ gives $\E_\mu[\mathcal G_\psi]=1$.
\end{proof}

\section{Equivalence}\label{sec:concl}

\begin{proof}[Proof of Theorem~\ref{thm:main}]
By definition of pushforward, for every bounded measurable $F$,
\[
  \int F\,\dd(T_\psi^*\mu)=\int F(T_\psi(\phi))\,\dd\mu(\phi)=\int F(\phi+\psi)\,\dd\mu(\phi).
\]
Replacing $\psi$ by $-\psi$ throughout Sections~\ref{sec:density}--\ref{sec:identity}
is legitimate because $-\psi$ is again smooth and not identically zero; equivalently,
apply Proposition~\ref{prop:limit} with the test function $\widetilde F(\phi)=F(\phi+\psi)$
which is bounded continuous on $\mathcal C^{-1/2-\kappa}$. Since
$\widetilde F(\phi-\psi)=F(\phi)$, identity \eqref{eq:CVlim} reads
\begin{equation}\label{eq:pushdens}
  \int F\,\dd\mu=\int F(\phi-\psi)\,\dd\mu(\phi)\Big|_{F\to\widetilde F}
  \ \Longrightarrow\
  \int F(\phi)\,\dd\mu(\phi)=\int F(\phi+\psi)\,\mathcal G_\psi(\phi)\,\dd\mu(\phi).
\end{equation}
Equivalently, applying \eqref{eq:CVlim} directly with the shift $\psi$ replaced by
$-\psi$ and density $\mathcal G_{-\psi}$ (the version of \eqref{eq:Phidef} with $\psi$
replaced by $-\psi$), we obtain
\begin{equation}\label{eq:final}
  \int F(\phi+\psi)\,\dd\mu(\phi)=\int F(\phi)\,\mathcal G_{-\psi}(\phi)\,\dd\mu(\phi),
\end{equation}
that is,
\begin{equation}\label{eq:RN}
  \frac{\dd(T_\psi^*\mu)}{\dd\mu}=\mathcal G_{-\psi}\quad\mu\text{-a.s.}
\end{equation}
By Proposition~\ref{prop:density} (applied to $-\psi$), $\mathcal G_{-\psi}>0$ everywhere
and $\mathcal G_{-\psi}\in L^1(\mu)$ with $\E_\mu[\mathcal G_{-\psi}]=1$
(Proposition~\ref{prop:limit}). Therefore $T_\psi^*\mu$ is absolutely continuous with
respect to $\mu$ with a strictly positive density.

Conversely, $T_\psi$ is a bijection of $\mathcal D'(\T)$ with inverse $T_{-\psi}$, and
$(T_\psi)^{-1}=T_{-\psi}$, so $\mu=T_{-\psi}^*(T_\psi^*\mu)$. Applying \eqref{eq:final}
with $\psi$ replaced by $-\psi$ gives $T_{-\psi}^*\mu\ll\mu$ with density
$\mathcal G_{\psi}>0$. Composing, $\mu\ll T_\psi^*\mu$ with strictly positive density
$\dd\mu/\dd(T_\psi^*\mu)=1/\mathcal G_{-\psi}\circ$\,(suitable shift), which is positive
and finite $\mu$-a.s.

Since both Radon--Nikodym derivatives \eqref{eq:RN} and its reciprocal are strictly
positive and finite $\mu$-a.s., $\mu$ and $T_\psi^*\mu$ are mutually absolutely
continuous, i.e.\ they have the same null sets. This is the desired equivalence, and the
answer to the problem is \textbf{yes}.
\end{proof}

\begin{remark}
The strict positivity in \eqref{eq:RN} is the reason equivalence (two-sided absolute
continuity) holds and not merely one-sided absolute continuity: a Radon--Nikodym density
that is $\mu$-a.s.\ strictly positive forces the same null sets. The mechanism is
classical Cameron--Martin quasi-invariance of the underlying Gaussian field $\nu$ — valid
precisely because $\psi$ lies in the Cameron--Martin space $H^1(\T)$, which every smooth
function does — promoted to $\mu$ through the renormalized change of the interaction
potential \eqref{eq:Vdiff}. The renormalization counterterms reorganize exactly into the
Wick powers $\Wick{\phi^3},\Wick{\phi^2}$, whose finiteness and exponential integrability
(Lemmas~\ref{lem:wick}--\ref{lem:tightabs}) keep the density \eqref{eq:G} a bona fide
$L^1(\mu)$ random variable despite the singularity of $\mu$ with respect to $\nu$.
\end{remark}

\end{document}
